% Part 1, Problem 13: Continued fraction for sqrt{2} via recursive sequence
% Hard

\begin{problem}[Continued Fraction for 1 + sqrt(2)]
Let $a = 1 + \sqrt{2}$.
\begin{enumerate}[label=(\roman*)]
    \item Show that $a^2 - 2a - 1 = 0$.
    \item Define $x_1 = 2$ and $x_{n+1} = 2 + \frac{1}{x_n}$ for $n \geq 1$. Assuming $x_n \to L$ for some positive limit $L$, prove that $L = a$.
    \item Show that $1 + \frac{1}{a} = \sqrt{2}$.
    \item Hence deduce that
    \[
    \sqrt{2} = 1 + \frac{1}{2 + \frac{1}{2 + \frac{1}{2 + \dots}}}.
    \]
\end{enumerate}
\end{problem}

\begin{hint}
Attempt the proof independently first. Focus on the key theorem, algebraic transformation, or contradiction setup that links the hypothesis to the target conclusion.
\end{hint}

\begin{solution}
\textbf{(i)} Since $a = 1+\sqrt{2}$, we have $a-1=\sqrt{2}$. Squaring:
\[
(a-1)^2=2 \implies a^2-2a+1=2 \implies a^2-2a-1=0.
\]

\textbf{(ii)} If $x_n \to L$, then also $x_{n+1} \to L$. Taking limits in
\[
x_{n+1}=2+\frac{1}{x_n}
\]
gives
\[
L=2+\frac{1}{L}.
\]
So
\[
L^2-2L-1=0.
\]
By part (i), $a$ satisfies the same quadratic. The roots are $1\pm\sqrt{2}$, and since $L>0$, we get
\[
L=1+\sqrt{2}=a.
\]

\textbf{(iii)} Using $a=1+\sqrt{2}$,
\[
1+\frac{1}{a}=1+\frac{1}{1+\sqrt{2}}
=1+\frac{1-\sqrt{2}}{(1+\sqrt{2})(1-\sqrt{2})}
=1-(1-\sqrt{2})=\sqrt{2}.
\]

\textbf{(iv)} The recursion generates
\[
x_1=2,\quad x_2=2+\frac{1}{2},\quad x_3=2+\frac{1}{2+\frac{1}{2}},\ \dots
\]
Hence its limit is exactly
\[
2+\frac{1}{2+\frac{1}{2+\dots}}.
\]
By part (ii), this limit is $a$, so part (iii) gives
\[
\sqrt{2}=1+\frac{1}{a}
=1+\frac{1}{2+\frac{1}{2+\frac{1}{2+\dots}}}.
\]

\hfill $\square$
\end{solution}

\begin{takeaways}
\begin{itemize}
\item \textbf{Limit Method:} Infinite continued fractions are handled by defining finite approximations and taking a limit
\item \textbf{Fixed Point Idea:} A convergent recursive sequence often leads to an equation satisfied by its limit
\item \textbf{Quadratic Link:} The continued fraction for $\sqrt{2}$ comes from solving a simple quadratic
\item \textbf{Representation:} This gives a classic infinite continued fraction expansion for $\sqrt{2}$
\end{itemize}
\end{takeaways}
